\section{Countermeasures}
\label{sec:countermeasures}

The controls fall into three groups: those that act before signing, those that act at the customer, and those that act after installation.

Before signing, the useful controls address the build itself. Separating build infrastructure from office networks, restricting outbound traffic, using short-lived agents, and limiting who can change build jobs all reduce the opportunity SUNSPOT used. CISA recommends protecting code from unauthorised access and using properly configured build processes~\cite{cisa-supply-chain}. Comparing a released artifact with an independently produced reference, or rebuilding from source and comparing hashes, targets the substitution directly. Neither proves the code is benign; a reproducible malicious build is still malicious.

At the customer, the signature is the wrong question. Least privilege, restricted update permissions, and a baseline of expected vendor traffic limit what an installed component can do. SBOMs and supplier assessments shorten the time to identify affected components but do not prevent the build compromise~\cite{cisa-supply-chain}.

After installation, detection depends on behaviour. DNS logging and anomaly analysis can show an unexpected check-in, and network baselines can separate vendor telemetry from a lookalike. EDR has little to work with while a signed component behaves normally, but becomes useful when the operator acts. CISA's Orion guidance covers reviewing stored traffic for new external DNS domains and comparing a product's expected external communications with what is observed~\cite{cisa-aa20-352a,cisa-supply-chain}.

Table~\ref{tab:countermeasures} rates each control for expected effectiveness and deployment difficulty. The ratings are the authors' assessment for a typical enterprise, not measurements from the lab.

\begin{table}[H]
\centering
\footnotesize
\caption{Countermeasures by stage, with expected effectiveness and deployment difficulty.}
\label{tab:countermeasures}
\begin{tabular}{@{}>{\RaggedRight}p{3.5cm}>{\RaggedRight}p{3.3cm}>{\RaggedRight}p{2.9cm}>{\RaggedRight}p{4.9cm}@{}}
\toprule
\textbf{Countermeasure} & \textbf{Effectiveness} & \textbf{Difficulty} & \textbf{Reason} \\
\midrule
Build pipeline isolation & High & High & Rebuilds infrastructure, network controls, and developer workflow \\
Reproducible builds & High for detecting tampering & Very high & Pinned dependencies, deterministic toolchains, controlled timestamps \\
Independent artifact verification & High & High & Needs a separate verification path, trusted records, and a release gate \\
HSM-protected signing keys & Low for this attack & Medium & Protects the signing step, which the attackers did not need to break \\
DNS anomaly detection & High for visibility & Medium & Needs logging and tuning; a fixed-domain lab is not a DGA test \\
Network behaviour analysis & Medium to high & Medium to high & Needs a baseline before deviations mean anything \\
EDR & Medium & Low to medium & Limited against a signed passive component \\
MFA and least privilege & Medium & Low & Limits movement, not the initial update \\
SBOM and vendor risk management & Low to medium now & Medium & Gives visibility rather than prevention \\
\bottomrule
\end{tabular}
\end{table}

The controls with the most effect are also the hardest to deploy, which is the practical problem with supply-chain defence. A customer cannot verify a build it did not perform, and a supplier cannot prove a build was clean without records the attacker could not alter.
